
    %%%%%%%%%%%%%%%%%%%%%%%%%%%%%%%%%%%%%%%%%%%%%%%%%%%%%%%%%%%%%%%%%%%%%%%%%%%%%%%%
    %%% Classe do documento
    \documentclass[12pt, addpoints]{exam}
    \UseRawInputEncoding

    %%%%%%%%%%%%%%%%%%%%%%%%%%%%%%%%%%%%%%%%%%%%%%%%%%%%%%%%%%%%%%%%%%%%%%%%%%%%%%%%
    % preambulo.tex
    %
    % Arquivo de configuração de packages para uso o LaTeX, conforme minhas
    % preferências e modelos pessoais.
    %
    % Para maiores informações, visite:
    %    https://github.com/abrantesasf/latex
    %
    % NÃO ALTERE SE NÃO SOUBER O QUE ESTÁ FAZENDO!
    %%%%%%%%%%%%%%%%%%%%%%%%%%%%%%%%%%%%%%%%%%%%%%%%%%%%%%%%%%%%%%%%%%%%%%%%%%%%%%%%


    %%%%%%%%%%%%%%%%%%%%%%%%%%%%%%%%%%%%%%%%%%%%%%%%%%%%%%%%%%%%%%%%%%%%%%%%%%%%%%%%
    %%% Estruturas de controle:
    \input{/app/static/utils/controles}

    %%%%%%%%%%%%%%%%%%%%%%%%%%%%%%%%%%%%%%%%%%%%%%%%%%%%%%%%%%%%%%%%%%%%%%%%%%%%%%%%
    %%% Compilação condicional em PDF:
    \input{/app/static/utils/pdf}

    %%%%%%%%%%%%%%%%%%%%%%%%%%%%%%%%%%%%%%%%%%%%%%%%%%%%%%%%%%%%%%%%%%%%%%%%%%%%%%%%
    %%% Configurações de layout da página:
    \input{/app/static/utils/layout}

    %%%%%%%%%%%%%%%%%%%%%%%%%%%%%%%%%%%%%%%%%%%%%%%%%%%%%%%%%%%%%%%%%%%%%%%%%%%%%%%%
    %%% Configurações para autores:
    \input{/app/static/utils/autores}

    %%%%%%%%%%%%%%%%%%%%%%%%%%%%%%%%%%%%%%%%%%%%%%%%%%%%%%%%%%%%%%%%%%%%%%%%%%%%%%%%
    %%% Cabeçalho e rodapé:
    %%% Atenção! É MUITO DIFÍCIL ajustar apropriadamente os running headers
    %%% e running footers da primeira vez. Você pode confiar no LaTeX e deixar que
    %%% ele faça o ajuste automaticamente ou pode quebrar a cabeça e
    %%% tentar melhorar algo que o LaTeX já faz muito bem, mesmo que não fique
    %%% exatamente do jeito que você quer. O que você prefere? Se quiser ajustar
    %%% manualmente, descomente a linha abaixo e faça as configurações no arquivo
    %%% "utils/headings.tex".
    %\input{/app/static/utils/headings}

    %%%%%%%%%%%%%%%%%%%%%%%%%%%%%%%%%%%%%%%%%%%%%%%%%%%%%%%%%%%%%%%%%%%%%%%%%%%%%%%%
    %%% Configuração para as fontes do documento:
    %%% Se você quiser usar fontes próprias ou do sistema, precisará ajustar as
    %%% configurações no arquivo "utils/fontes.tex".
    %%% ATENÇÃO: por padrão eu utilizo XeLaTeX com fontes proprietárias projetadas
    %%% por Matthew Butterick, adquiridas em https://mbtype.com, que não podem ser
    %%% distribuídas devido à licença de utilização. Se você usar XeLaTeX, você
    %%% DEVERÁ OBRIGATORIAMENTE ajustar as fontes no arquivo "utils/fontes.tex".
    \input{/app/static/utils/fontes}

    %%%%%%%%%%%%%%%%%%%%%%%%%%%%%%%%%%%%%%%%%%%%%%%%%%%%%%%%%%%%%%%%%%%%%%%%%%%%%%%%
    %%% Sumário:
    \input{/app/static/utils/toc}

    %%%%%%%%%%%%%%%%%%%%%%%%%%%%%%%%%%%%%%%%%%%%%%%%%%%%%%%%%%%%%%%%%%%%%%%%%%%%%%%%
    %%% Matemática:
    \input{/app/static/utils/matematica}

    %%%%%%%%%%%%%%%%%%%%%%%%%%%%%%%%%%%%%%%%%%%%%%%%%%%%%%%%%%%%%%%%%%%%%%%%%%%%%%%%
    %%% Notas de rodapé:
    \input{/app/static/utils/footnotes}

    %%%%%%%%%%%%%%%%%%%%%%%%%%%%%%%%%%%%%%%%%%%%%%%%%%%%%%%%%%%%%%%%%%%%%%%%%%%%%%%%
    %%% Referências cruzadas, links, citações:
    \input{/app/static/utils/crossrefs}

    %%%%%%%%%%%%%%%%%%%%%%%%%%%%%%%%%%%%%%%%%%%%%%%%%%%%%%%%%%%%%%%%%%%%%%%%%%%%%%%%
    %%% Configurações para os hiperlinks em PDF:
    \input{/app/static/utils/pdfconfig}

    %%%%%%%%%%%%%%%%%%%%%%%%%%%%%%%%%%%%%%%%%%%%%%%%%%%%%%%%%%%%%%%%%%%%%%%%%%%%%%%%
    %%% Glossário e índice remissivo:
    \input{/app/static/utils/glossindex}

    %%%%%%%%%%%%%%%%%%%%%%%%%%%%%%%%%%%%%%%%%%%%%%%%%%%%%%%%%%%%%%%%%%%%%%%%%%%%%%%%
    %%% Referências bibliográficas:
    \input{/app/static/utils/bibliografia}

    %%%%%%%%%%%%%%%%%%%%%%%%%%%%%%%%%%%%%%%%%%%%%%%%%%%%%%%%%%%%%%%%%%%%%%%%%%%%%%%%
    %%% Cores:
    \input{/app/static/utils/cores}

    %%%%%%%%%%%%%%%%%%%%%%%%%%%%%%%%%%%%%%%%%%%%%%%%%%%%%%%%%%%%%%%%%%%%%%%%%%%%%%%%
    %%% Computação:
    \input{/app/static/utils/computacao}

    %%%%%%%%%%%%%%%%%%%%%%%%%%%%%%%%%%%%%%%%%%%%%%%%%%%%%%%%%%%%%%%%%%%%%%%%%%%%%%%%
    %%% Gráficos:
    \input{/app/static/utils/graficos}

    %%%%%%%%%%%%%%%%%%%%%%%%%%%%%%%%%%%%%%%%%%%%%%%%%%%%%%%%%%%%%%%%%%%%%%%%%%%%%%%%
    %%% Ícones:
    \input{/app/static/utils/icones}

    %%%%%%%%%%%%%%%%%%%%%%%%%%%%%%%%%%%%%%%%%%%%%%%%%%%%%%%%%%%%%%%%%%%%%%%%%%%%%%%%
    %%% Blocos:
    \input{/app/static/utils/blocos}

    %%%%%%%%%%%%%%%%%%%%%%%%%%%%%%%%%%%%%%%%%%%%%%%%%%%%%%%%%%%%%%%%%%%%%%%%%%%%%%%%
    %%% Boxes coloridos:
    \input{/app/static/utils/colorbox}

    %%%%%%%%%%%%%%%%%%%%%%%%%%%%%%%%%%%%%%%%%%%%%%%%%%%%%%%%%%%%%%%%%%%%%%%%%%%%%%%%
    %%% Tabelas:
    \input{/app/static/utils/tabelas}

    %%%%%%%%%%%%%%%%%%%%%%%%%%%%%%%%%%%%%%%%%%%%%%%%%%%%%%%%%%%%%%%%%%%%%%%%%%%%%%%%
    %%% Ambientes de listas:
    \input{/app/static/utils/listas}

    %%%%%%%%%%%%%%%%%%%%%%%%%%%%%%%%%%%%%%%%%%%%%%%%%%%%%%%%%%%%%%%%%%%%%%%%%%%%%%%%
    %%% Ambientes floats e similares:
    \input{/app/static/utils/floats}

    %%%%%%%%%%%%%%%%%%%%%%%%%%%%%%%%%%%%%%%%%%%%%%%%%%%%%%%%%%%%%%%%%%%%%%%%%%%%%%%%
    %%% Ajustes para a classe EXAM:
    \input{/app/static/utils/exam}

    %%%%%%%%%%%%%%%%%%%%%%%%%%%%%%%%%%%%%%%%%%%%%%%%%%%%%%%%%%%%%%%%%%%%%%%%%%%%%%%%
    %%% Ajustes para textos:
    \input{/app/static/utils/texto}

    %%%%%%%%%%%%%%%%%%%%%%%%%%%%%%%%%%%%%%%%%%%%%%%%%%%%%%%%%%%%%%%%%%%%%%%%%%%%%%%%
    %%% Ambientes, aliases, comandos e customizações em geral para serem
    %%% utilizadas nos documentos. Defina aqui qualquer customização específica
    %%% para seus documentos!
    \input{/app/static/utils/customizacoes}

    %%%%%%%%%%%%%%%%%%%%%%%%%%%%%%%%%%%%%%%%%%%%%%%%%%%%%%%%%%%%%%%%%%%%%%%%%%%%%%%%
    %%% Ajuste do layout, espaçamento de linhas e etc.:
    \geometry{a4paper, portrait, left=2.0cm, right=2.0cm, top=2.0cm, bottom=2.0cm}
    \singlespacing

    %%%%%%%%%%%%%%%%%%%%%%%%%%%%%%%%%%%%%%%%%%%%%%%%%%%%%%%%%%%%%%%%%%%%%%%%%%%%%%%%
    %%% Configurações de cabeçalho e rodapé (não use fancyhdr pois ocorre conflito):
    \pagestyle{headandfoot}
    \runningheadrule
    \firstpageheader{}{}{}
    \runningheader{Sem disciplina}{}{Avaliação}
    \firstpagefooter{}{}{Página \thepage\ de \numpages}
    \runningfooter{}{}{Página \thepage\ de \numpages}
    \runningfootrule

    %%%%%%%%%%%%%%%%%%%%%%%%%%%%%%%%%%%%%%%%%%%%%%%%%%%%%%%%%%%%%%%%%%%%%%%%%%%%%%%%
    %%% Configurações para as propriedades do PDF:
    \hypersetup{
    hidelinks,           % Comente para web, descomente para impressão
    colorlinks=true,      % True para web, False para impressão
    pdftitle=jjj: Avaliação,
    pdfauthor=be@qe.com,
    pdfsubject=Sem disciplina,
    pdfkeywords=['Sem tag'],
    pdfinfo={
        CreationDate={}, % Ex.: D:AAAAMMDDHH24MISS
        ModDate={}       % Ex.: D:AAAAMMDDHH24MISS
    }
    }
    \input{/app/static/utils/pdfconfig}

    %%%%%%%%%%%%%%%%%%%%%%%%%%%%%%%%%%%%%%%%%%%%%%%%%%%%%%%%%%%%%%%%%%%%%%%%%%%%%%%%
    %%% Começa o documento
    \begin{document}

    %%%%%%%%%%%%%%%%%%%%%%%%%%%%%%%%%%%%%%%%%%%%%%%%%%%%%%%%%%%%%%%%%%%%%%%%%%%%%%%%
    %%% Folha de instruções padronizadas da UVV para a prova
    \pdfbookmark[1]{Instruções}{instrucoes}
    \begin{figure}[H]
        \begin{center}
            \includegraphics[scale=0.3]{{/app/static/imgs/logo-uvv.png}}
        \end{center}	
    \end{figure}

    \vspace{-1cm}

    \begin{table}[H]
        \centering
        \begin{tabular}{|llll|}
            \hline
            \multicolumn{2}{|l|}{\textbf{Disciplina:} Sem disciplina} & 
            \multicolumn{1}{l|}{\multirow{3}{*}{\textbf{\makecell{Nota:\\ \,\\ \,}}}} & 
            \multirow{3}{*}{\textbf{\makecell{Coordenador:\\ \,\\ \,}}} \\ 
            \cline{1-2}
            \multicolumn{2}{|l|}{\textbf{Professor:} be@qe.com \hspace{4.72cm} } & 
            \multicolumn{1}{l|}{} & \\ 
            \cline{1-2}
            \multicolumn{2}{|l|}{\textbf{Aluno:}} & 
            \multicolumn{1}{l|}{} & \\ 
            \hline
            \multicolumn{1}{|l|}{\textbf{Turma:} \hspace{6.3cm} } & 
            \multicolumn{1}{l|}{\textbf{Semestre:}} & 
            \multicolumn{2}{l|}{\textbf{Valor:} 10 pontos} \\ 
            \hline
            \multicolumn{1}{|l|}{\textbf{Data:}} & 
            \multicolumn{3}{l|}{\textbf{Avaliação:}} \\ 
            \hline
        \end{tabular}
    \end{table}

    \begin{center}
        \fbox{\setlength\fboxsep{0.3cm} \fbox{\parbox{15.4cm}{
            \begin{center}
                \textbf{INSTRUÇÕES PARA A REALIZAÇÃO DESTA AVALIAÇÃO:}
            \end{center}
            \begin{itemize}[leftmargin=*]
                \item Esta é a primeira avaliação da disciplina Sem disciplina, 
                    e engloba o conteúdo referente Sem tag.
                \item Esta avaliação contém 40 questões objetivas, \textbf{todas obrigatórias}.
                \item \textbf{Leia atentamente cada questão!} As questões objetivas terão uma,
                    e apenas uma única, resposta a ser marcada.
                \item \textbf{Desligue o celular} antes de começar. A avaliação é
                    \textbf{individual} e \textbf{sem consulta}.
                \item As questões podem ser respondidas, na prova, com lápis ou caneta. Ao final
                    da prova \textbf{{VOCÊ DEVERÁ PREENCHER O GABARITO COM CANETA
                    DE COR PRETA}} \textbf{\underline{OU AZUL}} para que seja feita a correção
                    automatizada através da leitura do padrão de respostas marcado no
                    gabarito.
                \item \textbf{CUIDADO ao preencher o gabarito} pois eles são \textbf{INDIVIDUAIS
                    e NOMEADOS} (cada estudante tem um gabarito próprio específico, com um
                    código de identificação único). \textbf{Questões rasuradas são anuladas}
                    pelo software de leitura do gabarito.
                \item Ao final da prova, \textbf{DEVOLVA A PROVA E O GABARITO} para o professor.
                \item Siga todas as normas de \textbf{Integridade Acadêmica} da disciplina.
                    Alunos que forem flagrados com qualquer espécie de ``cola'' ou trocando
                    informações com outros alunos terão suas avaliações recolhidas, as notas
                    zeradas, e a situação será encaminhada para a coordenação para a aplicação
                    das penalidades previstas pela universidade.
                \item Com 90 minutos de avaliação você terá tempo de até 2,min por questão,
                    em média.
                \item Boa avaliação!
            \end{itemize}
            \vspace{2cm}
        }}}
    \end{center}
    \newpage

    %%%%%%%%%%%%%%%%%%%%%%%%%%%%%%%%%%%%%%%%%%%%%%%%%%%%%%%%%%%%%%%%%%%%%%%%%%%%%%%%
    %%% Inicia o bloco de questões:
    \begin{questions}

    \newpage
    \question
Considere a seguinte tabela em uma base de dados relacional (chave primária sublinhada):\\
Tabela1 (CodAluno, CodDisciplina, AnoSemestre, NomeAluno, NomeDisciplina, CodNota, DescricaoNota)\\
Considere as seguintes dependências funcionais:
\begin{itemize}
    \item CodAluno \(\rightarrow\) NomeAluno
    \item CodDisciplina \(\rightarrow\) NomeDisciplina
    \item (CodAluno, CodDisciplina, AnoSemestre) \(\rightarrow\) CodNota
    \item (CodAluno, CodDisciplina, AnoSemestre) \(\rightarrow\) DescricaoNota
    \item CodNota \(\rightarrow\) DescricaoNota
\end{itemize}

Considerando as formas normais, qual das afirmativas abaixo se aplica:
\begin{choices}
\choice A tabela não está na primeira forma normal.
\choice A tabela encontra-se na segunda forma normal, mas não na terceira forma normal.
\choice A tabela encontra-se na terceira forma normal, mas não na quarta forma normal.
\choice A tabela está na quarta forma normal.
\choice A tabela encontra-se na primeira forma normal, mas não na segunda forma normal.
\end{choices}

\question
\begin{figure}[H]
    \centering
    \includegraphics[width=0.5\linewidth]{/app/static/imgs/defaul.png}
    \caption{Enter Caption}
\end{figure}
Se \( O = (0, 0, 0) \); \( A = (2, 4, 1) \); \( B = (3, 1, 1) \) e \( C = (1, 3, 5) \) então o volume do sólido acima é

\begin{choices}
\choice 35
\choice 44
\choice 21
\choice 30
\choice 35/2
\end{choices}

\question
Um banco de dados OLAP (Online Analytical Processing --- Processamento Analítico
Online) é um banco de dados organizado de forma a dar suporte ao Business
Intelligence (BI --- Inteligência de Negócio) que, de forma simplificada, é o
processo de analisar dados para obter informações que podem ser utilizadas na
tomada de decisões comerciais e estratégicas. Em relação aos bancos de dados
OLAP, marque a afirmação correta:
\begin{choices}
\choice São otimizados para consulta e relatórios de rotina de dados.
\choice São preparados para responder perguntas do tipo ``Quais alunos estão matriculados na turma de Bancos de Dados I neste semestre?'
\choice São otimizados para o processamento de grande volume de dados históricos e métricas de negócio para a tomada de decisões estratégicas, permitindo e facilitando a análise de dados para a geração de informações e conhecimentos para a tomada de decisões.
\choice Geralmente os dados de origem do OLTP são de bancos de dados OLAP.
\choice São otimizados para o processamento de grande volume de transações diárias, o que permite a análise e a tomada de decisões em tempo real, favorecendo assim a administração empresarial.
\end{choices}

\question
O contador da figura abaixo é:
\begin{figure}
    \centering
    \includegraphics[width=0.5\linewidth]{/app/static/imgs/defaul.png}
    \label{fig:enter-label}
\end{figure}
\begin{choices}
\choice anisócrono
\choice assíncrono
\choice auto-sincronizado
\choice síncrono
\choice isócrono
\end{choices}

\question
Sistemas de processamento de transações, tais como sistemas de reservas aéreas, devem prover um mecanismo que garanta que cada transação não é afetada por outras transações que possam estar ocorrendo ao mesmo tempo. Transações de duas fases obedecem a um protocolo que garante essa atomicidade. Em transações de duas fases:

\begin{choices}
\choice Todas as operações de leitura ocorrem antes da primeira operação de escrita.
\choice Verifica-se a disponibilidade de todas as travas antes de executar qualquer ação de travamento.
\choice Uma trava compartilhada sobre um objeto deve ser obtida antes de uma trava exclusiva sobre o objeto ser obtida.
\choice Todas as ações de travamento (\textit{lock}) ocorrem antes da primeira ação de destravamento.
\choice Qualquer objeto corretamente travado deve ser destravado antes que outro objeto possa ser travado.
\end{choices}

\question
Na linguagem SQL, qual das seguintes palavras-chave é usada para modificar os
dados em uma tabela existente?

\begin{choices}
\choice CREATE
\choice SELECT
\choice Nenhuma das alternativas acima
\choice INSERT
\choice UPDATE
\end{choices}

\question
Considere as seguintes afirmações sobre mecanismos de inferência em sistemas baseados em regras.
\begin{enumerate}[label=\Roman*.]
    \item O encadeamento regressivo tem pouca utilidade prática, pois deve partir do possível resultado.
    \item O encadeamento progressivo tanto pode ser em amplitude quanto em profundidade.
    \item Podem trabalhar com informações incertas ou incompletas. 
\end{enumerate}
São corretas:
\begin{choices}
\choice Apenas III
\choice Apenas I e II
\choice Apenas I e III
\choice I, II e III
\choice Apenas II e III
\end{choices}

\question
Considere a seguinte tabela para uma base de dados relacional:\\
Empregado (CodEmp, NomeEmp, CodDepto)\\
Considere que esta tabela tem um índice na forma de uma árvore B sobre as colunas (CodEmp, CodDepto), nesta ordem.
Quanto a este índice, considere as seguintes afirmativas:
\begin{enumerate}
    \item Este índice pode ser usado pelo SGBD relacional para acelerar uma consulta na qual são fornecidos os valores de CodEmp e CodDepto.
    \item Este índice pode ser usado pelo SGBD relacional para acelerar uma consulta na qual é fornecido um valor de CodEmp.
    \item Este índice não é adequado para ser usado pelo SGBD relacional para acelerar uma consulta na qual é fornecido um valor de CodDepto.
    \item O algoritmo que faz inserções e remoções de entradas do índice tem por objetivo garantir que o índice fique organizado de tal forma que o acesso a cada nodo da árvore implique em número de acessos semelhantes.
    \item O índice por árvore-B não é adequado para tabelas que sofrem grande número de inclusões e exclusões, pois exige reorganizações frequentes.
\end{enumerate}


Quanto a estas afirmativas pode se dizer que:
\begin{choices}
\choice Nenhuma das afirmativas está correta
\choice Apenas as afirmativas 1), 2) e 5) estão corretas
\choice Apenas as afirmativas 1), 2) e 4) estão corretas
\choice Todas afirmativas estão corretas
\choice Apenas as afirmativas 1), 2), 3) e 4) estão corretas
\end{choices}

\question
Quando trabalhando com sistemas baseados em trocas de mensagens, temporizações (\textit{time-outs}) são utilizadas para:

\begin{choices}
\choice Limitar o tamanho de uma mensagem transmitida.
\choice Limitar o número de retransmissões de uma mensagem.
\choice Arbitrar que uma mensagem transmitida foi perdida.
\choice Temporariamente suspender a transmissão de mensagens.
\choice Limitar o tempo para obter um recurso.
\end{choices}

\question
Em relação aos metadados de um banco de dados, assinale a afirmação correta:
\begin{choices}
\choice Armazenam informações sobre os dados dos usuários, ou seja, são dados de rotina dos usuários, aqueles dados que os usuários precisam para trabalhar no seu dia a dia.
\choice São armazenados no catálogo (ou dicionário de dados) do sistema de gerenciamento de bancos de dados, que fica nas mesmas tabelas de dados dos usuários, permitindo assim que os metadados estejam sempre atualizados.
\choice Como são sempre atualizados de forma automática pelo sistema de gerenciamento de bancos de dados, estão sempre prontos para nos dar uma ``foto' da estrutura do banco de dados em um momento no tempo.
\choice São atualizados pelo SGBD, sempre que ocorre alguma operação de inserção, atualização ou remoção de dados.
\choice Armazenam informações como os dados de cadastro dos usuários, os valores de compra em uma ordem de compra, as turmas e disciplinas que os professores dão aula, e dados de recursos humanos.
\end{choices}

\question
Em relação ao teste de software, qual das afirmações a seguir é INCORRETA: 
\begin{choices}
\choice Um bom caso de teste é aquele que tem uma elevada probabilidade de revelar um erro ainda não descoberto. 
\choice O processo de depuração é a parte mais imprevisível do processo de teste pois um erro pode demorar uma hora, um dia ou um mês para ser diagnosticado e corrigido. 
\choice A atividade de teste é o processo de executar um programa com a intenção de demonstrar a ausência de erros. 
\choice Os dados compilados quando a atividade de teste é levada a efeito proporcionam uma boa indicação da confiabilidade do software e alguma indicação da qualidade do software como um todo. 
\choice Um teste bem sucedido é aquele que revela um erro ainda não descoberto. 
\end{choices}

\question
A medida da interconexão entre os módulos de uma estrutura de software é denominada e que também é usada em projetos orientados a objetos é: 
\begin{choices}
\choice unidade funcional
\choice coesão
\choice ocultamento da informação
\choice abstração procedimental
\choice acoplamento
\end{choices}

\question
Seja \( f : \mathbb{R} \to \mathbb{R} \) derivável. Se existem \( a, b \in \mathbb{R} \) tal que \( f(a)f(b) < 0 \) e \( f(x) \neq 0 \) para todo \( x \in (a, b) \), podemos afirmar que no intervalo \( (a, b) \) a equação \( f(x) = 0 \) tem:
\begin{choices}
\choice duas raízes reais
\choice uma única raiz real
\choice uma raiz imaginária
\choice somente raízes imaginárias
\choice nenhuma raiz real
\end{choices}

\question
Em um banco de dados relacional, o que é uma chave primária?

\begin{choices}
\choice É denominada por FK
\choice Nenhuma das respostas acima
\choice É uma validação que evita o cadastro de dados errados
\choice É um atributo (ou um conjunto de atributos) que permite identificar única e exclusivamente cada registro da tabela
\choice É qualquer coluna ou conjunto de colunas que podem servir como índice
\end{choices}

\question
Qual dos protocolos abaixo pode ser caracterizado como protocolo de roteamento do tipo estado de enlace?
\begin{choices}
\choice IGMP
\choice ICMP
\choice OSPF
\choice BGP-4
\choice RIP2
\end{choices}

\question
Qual o valor do atributo \( E.val \) após a análise da expressão "4 / 2 / 2" para o esquema de tradução a seguir?

\begin{itemize}
    \item \( E \rightarrow T / E1 \ \{ E.val = T.val / E1.val \} \)
    \item \( E \rightarrow T \ \{ E.val = T.val \} \)
    \item \( T \rightarrow \text{digito} \ \{ T.val = \text{val(digito)} \} \)
\end{itemize}

\begin{choices}
\choice 8
\choice 2
\choice 1
\choice 3
\choice 4
\end{choices}

\question
Em SQL, qual comando é utilizado para combinar linhas de duas ou mais tabelas?

\begin{choices}
\choice MERGE
\choice JOIN
\choice CONNECT
\choice LINK
\choice COMBINE
\end{choices}

\question
O processo de visualização de objetos 3D envolve uma série de passos desde a representação vetorial de um objeto até a exibição da imagem correspondente na tela do computador (pipeline 3D). Selecione a alternativa abaixo que reflete a ordem correta em que esses passos devem ocorrer. 
\begin{choices}
\choice Transformação de câmera, recorte 3D, projeção, mapeamento para coordenadas de tela, rasterização. 
\choice Projeção, transformação de câmera, recorte 3D, mapeamento para coordenadas de tela, rasterização. 
\choice Recorte 3D, transformação de câmera, rasterização, projeção, mapeamento para coordenadas de tela. 
\choice Transformação de câmera, mapeamento para coordenadas de tela, recorte 3D, rasterização, projeção. 
\choice Nenhuma das respostas acima está correta.
\end{choices}

\question
Na Álgebra Booleana:
\begin{choices}
\choice Os dígitos são octais, de 0 a 7
\choice Há dez valores motivados pelos dez dedos do ser humano
\choice Os dígitos são hexadecimais de 0 a 15
\choice Os dígitos são binários 0 e 1
\choice Os dígitos são alfanuméricos que podem ser representados por um byte
\end{choices}

\question
Considere as seguintes tabelas em uma base de dados relacional (chaves primárias sublinhadas):\\
Departamento (\underline{CodDepto}, NomeDepto)\\
Empregado (\underline{CodEmp}, NomeEmp, CodDepto)\\
Considere as seguintes restrições de integridade sobre esta base de dados relacional:\\
- Empregado.CodDepto é sempre diferente de NULL\\
- Empregado.CodDepto é chave estrangeira da tabela Departamento com cláusulas ON DELETE RESTRICT e ON UPDATE RESTRICT\\
Qual das seguintes validações não é especificada por estas restrições de integridade:
\begin{choices}
\choice Sempre que uma nova linha for inserida em Departamento, deve ser garantido que o valor de Departamento.CodDepto aparece na coluna Empregado.CodDepto.
\choice Sempre que o valor de Departamento.CodDepto for alterado, deve ser garantido que não há uma linha com o antigo valor de Departamento.CodDepto na coluna Empregado.CodDepto.
\choice Sempre que uma nova linha for inserida em Empregado, deve ser garantido que o valor de Empregado.CodDepto aparece na coluna Departamento.CodDepto.
\choice Sempre que o valor de Empregado.CodDepto for alterado, deve ser garantido que o novo valor de Empregado.CodDepto aparece em Departamento.CodDepto.
\choice Sempre que uma linha for excluída de Departamento, deve ser garantido que o valor de Departamento.CodDepto não aparece na coluna Empregado.CodDepto.
\end{choices}

\question
Sobre a UML, quais das seguintes afirmações são verdadeiras?

\begin{itemize}
    \item[I)] A UML é o método de desenvolvimento de software mais utilizado na atualidade.
    \item[II)] A UML é uma evolução das linguagens para especificação dos conceitos dos métodos de Booch, OMT e OOSE e também de outros métodos de especificação de requisitos de software orientados a objetos ou não.
    \item[III)] A UML é composta dos seguintes diagramas: Diagrama de Caso de Uso, Diagrama de Classes, Diagrama de Colaboração, Diagrama de Estados, entre outros.
    \item[IV)] Em UML pode-se representar tão somente relacionamentos de Agregação, Associação e Composição.
\end{itemize}
\begin{choices}
\choice Apenas as alternativas I, II e III.
\choice Apenas as alternativas II e III.
\choice Apenas as alternativas III e IV.
\choice Nenhuma delas.
\choice Todas as alternativas.
\end{choices}

\question
Engenharia de Software inclui um grande número de teorias, conceitos, modelos, técnicas e métodos. Analise as seguintes definições.
\begin{enumerate}[label=\Roman*.]
    \item O processo de inferir ou reconstruir um modelo de mais alto nível (projeto ou especificação) a partir de um documento de mais baixo nível (tipicamente um código fonte);
    \item Capacidade de modificação de um software (ou de um de seus componentes) após sua entrega ao cliente visando corrigir falhas, expandir a funcionalidade, modificar a performance ou outros atributos em resposta a novos requisitos do usuário ou mesmo ser adaptado a alguma mudança do ambiente de execução (plataforma, p.ex);
    \item Modelo estabelecido pelo Software Engineering Institute (SEI) que propõe níveis de competência organizacional relacionados à qualidade do processo de desenvolvimento de software;
\end{enumerate}
Estas definições correspondem respectivamente aos seguintes termos:
\begin{choices}
\choice engenharia reversa, manutenibilidade, Capability Maturity Model (CMM)
\choice reengenharia, manutenibilidade, Capability Maturity Model (CMM)
\choice engenharia reversa, reparabilidade, Team Software Process (TSP)
\choice reengenharia, evolutibilidade, Personal Software Process (PSP)
\choice engenharia reversa, reparabilidade, Team Software Process (TSP)
\end{choices}

\question
Seu chefe que fazer uma campanha de venda para produtos que custam entre 20 e 30
reais, mas está preocupado porque ouviu de um dos vendedores que diversos
produtos nessa faixa de preço estão com o estoque zerado. Para verificar a
situação seu chefe pediu para que você prepare um relatório que mostre os
produtos que custam entre 20 e 30 reais e que estão com o estoque zerado. Você
preparou o seguinte relatório para seu chefe:
\begin{tcolorbox}
 \begin{verbatim}1  SELECT l.nome AS loja
2       , p.nome AS produto
3       , p.preco_unitario AS preco
4       , e.quantidade AS qtd
5  FROM produtos p
6  INNER JOIN estoques e ON (e.produto_id = p.produto_id)
7  INNER JOIN lojas l ON (e.loja_id = l.loja_id)
8  WHERE (p.preco_unitario >= 20 OR p.preco_unitario <= 30)
9    AND e.quantidade = 0
10 ORDER BY loja, produto;
 \end{verbatim}
\end{tcolorbox}
O seu relatório:
\begin{choices}
\choice Nenhuma das alternativas anteriores.
\choice Está correto e trará exatamente, os produtos com o estoque zerado e na faixa de preço requisitada pelo seu chefe. Além disso ordena o resultado pelo nome da loja e pelo nome do produto.
\choice Está errado pois há um erro de sintaxe nas linhas 6 e 7.
\choice Está errado pois a tabela que deveria estar na cláusula FROM é a tabela ``lojas', não a tabela ``produtos'.
\choice Está errado pois vai trazer também produtos fora da faixa de preço que o seu chefe quer.
\end{choices}

\question
Um agente SNMP é um aplicativo que é executado:
\begin{choices}
\choice em "firewalls" com o objetivo de proteger acesso a rede
\choice em computadores denominados de gerentes
\choice a partir de um computador específico para monitorar a rede
\choice em um dispositivo de rede
\choice em roteadores com filtragem de pacotes com o objetivo de proteger acesso a rede
\end{choices}

\question
Professor Mac Sperto propôs o seguinte algoritmo de ordenação, chamado de Super Merge, similar ao \textit{merge sort}: divida o vetor em 4 partes do mesmo tamanho (ao invés de 2, como é feito no \textit{merge sort}). Ordene recursivamente cada uma das partes e depois intercale-as por um procedimento semelhante ao procedimento de intercalação do \textit{merge sort}. Qual das alternativas abaixo é verdadeira?
\begin{choices}
\choice Nenhuma das afirmações acima está correta
\choice Super Merge está correto, mas consome tempo \( O(\text{merge sort}) \)
\choice Super Merge não está correto. Não é possível ordenar quebrando o vetor em 4 partes
\choice Super Merge está correto, mas consome tempo maior que \( O(\text{merge sort}) \)
\choice Super Merge está correto, mas consome tempo menor que \( O(\text{merge sort}) \)
\end{choices}

\question
Uma característica de uma arquitetura RISC é:
\begin{choices}
\choice A arquitetura é constituída de milhares de processadores
\choice Possui um grande conjunto de instruções complexas
\choice O processador é formado por válvulas e transistores
\choice Possui um pequeno conjunto de instruções simples
\choice Uma arquitetura de alto risco pois o mercado de hardware evolui muito rapidamente
\end{choices}

\question
Para que serve a segmentação de um processador (\textit{pipelining})?
\begin{choices}
\choice reduzir o número de instruções estáticas nos programas
\choice simplificar o conjunto de instruções
\choice aumentar a velocidade do relógio
\choice permitir a execução de mais de uma instrução por ciclo de relógio
\choice simplificar a implementação do processador
\end{choices}

\question
Pode-se afirmar que o gráfico da função \(y = 2 + \frac{1}{x - 1}\) é o gráfico da função \(y = \frac{1}{x}\):
\begin{choices}
\choice transladado uma unidade para a direita e duas unidades para cima
\choice transladado uma unidade para a direita e duas unidades para baixo
\choice transladado uma unidade para a esquerda e duas unidades para cima
\choice nenhuma das anteriores.
\choice transladado uma unidade para a esquerda e duas unidades para baixo
\end{choices}

\question
Considere as seguintes afirmações sobre resolução de problemas em IA.
\begin{enumerate}[label=\Roman*.]
    \item A* é um conhecido algoritmo de busca heurística.
    \item O Minimax é um dos principais algoritmos para jogos de dois jogadores, como o xadrez.
    \item Busca em espaço de estados é uma das formas de resolução de problemas em IA.
\end{enumerate}
São corretas:
\begin{choices}
\choice Apenas III
\choice Apenas I e III
\choice I, II e III
\choice Apenas I e II
\choice Apenas II e III
\end{choices}

\question
Considere o circuito abaixo, implementado com duas portas \textbf{NAND}.

\begin{figure}[H]
    \centering
    \includegraphics[width=0.5\linewidth]{/app/static/imgs/defaul.png}
    \caption{Questão 23}
    \label{fig:enter-label}
\end{figure}

Qual das seguintes portas equivale a este circuito?
\begin{choices}
\choice \textbf{NOR}
\choice \textbf{AND}
\choice \textbf{XOR}
\choice \textbf{OR}
\choice \textbf{NOT}
\end{choices}

\question
Analise o diagrama relacional do banco de dados exibido na figura abaixo:
\begin{figure}[H]
  \centering
  \caption{Banco de dados de peças e fornecedores}
  \label{fig:pecas}
  %\fbox{
    \includegraphics[width=0.5\linewidth]{/app/static/imgs/defaul.png}
  %}\\
  %\footnotesize{Fonte: xxxx.}
\end{figure}
Podemos dizer que esse banco de dados ilustra:
\begin{choices}
\choice Um relacionamento com identificação N:N entre as tabelas ``pecas' e ``fornecedores', realizado através do relacionamento com cardinalidade identificada através da tabela ``fornecimentos'. 
\choice Um relacionamento com cardinalidade N:N entre as tabelas ``pecas' e ``fornecedores', realizado através de relacionamentos não identificados através da tabela ``fornecimentos'. 
\choice Um relacionamento com cardinalidade N:N entre as tabelas ``pecas' e ``fornecimentos', realizado através de relacionamentos não identificados com a tabela ``fornecedores'. 
\choice Um relacionamento com cardinalidade N:N entre as tabelas ``pecas' e ``fornecedores', realizado através de relacionamentos identificados através da tabela ``fornecimentos'. 
\choice Um relacionamento com cardinalidade 1:N entre as tabelas ``pecas' e ``fornecedores', indicando que uma peça pode ter sido fornecida por diversos fornecedores. 
\end{choices}

\question
Quanto ao TCP, é INCORRETO afirmar:
\begin{choices}
\choice Utiliza portas para permitir a comunicação entre processos localizados em dispositivos diferentes.
\choice Possui um campo de \textit{checksum} que valida as informações de seu cabeçalho, mas não valida as informações de \textit{payload} (campo de dados).
\choice É um protocolo do nível de transporte.
\choice Usa janelas deslizantes para implementar o controle de fluxo e erro.
\choice É um protocolo orientado a conexão.
\end{choices}

\question
A derivada da função \( f(x) = x^x \) é igual a:
\begin{choices}
\choice \( x^x (\ln(x) + 1) \)
\choice \( x^x \ln(x) \)
\choice \( x x^{x-1} \)
\choice \( x^x \)
\choice \( x^x (\ln(x) + x) \)
\end{choices}

\question
Quando Edgar Codd cricou o modelo relacional de dados,
também criou algumas regras para que a estrutura das tabelas fosse adequada às
operações de inserção, atualização, remoção e busca de dados. Ao avaliar um
banco de dados que armazena dados dos eleitores, dos partidos e dos candidatos
de uma eleição, você encontrou a situação ilustrada na tabela abaixo:
\begin{figure}[H]
  \centering
  \includegraphics[width=0.5\linewidth]{/app/static/imgs/defaul.png}
\end{figure}
Em relação às boas normas de projeto do modelo relacional, assinale a
alternativa que corretamente identifica um problema no projeto desta tabela:
\begin{choices}
\choice A chave primária para a tabela não está definida e, somente olhando as colunas, não é possível dizer se é a coluna 'numero' ou a coluna 'titulo'.
\choice A tabela está misturando dados de três entidades diferentes, eleitores, partidos e candidatos (o correto seria que cada tabela armazenasse dados de um único assunto ou entidade).
\choice Todas as afirmativas anteriores.
\choice A tabela contém um atributo multivalorado (o correto seria que cada tabela não tivesse atributos multivalorados).
\choice Apesar do telefone ser um atributo multivalorado, do jeito que a tabela está criada só é possível armazenar dois números de telefones, criando uma restrição ao armazenamento de telefones de contado de eleitores que tenham mais de dois telefones.
\end{choices}

\question
Uma prova de vestibular foi elaborada com 25 questões de múltipla escolha com 5 alternativas. O número de candidatos presentes à prova foi 63127. Considere a afirmação: Pelo menos 2 candidatos responderam de modo idêntico as \( k \) primeiras questões da prova. Qual é o maior valor de \( k \) para o qual podemos garantir que a afirmação é verdadeira.

\begin{choices}
\choice 10
\choice 9
\choice 7
\choice 8
\choice 6
\end{choices}

\question
Quais dos algoritmos de ordenação abaixo possuem tempo no pior caso e tempo médio de execução proporcional a \( O(n \log n) \).
\begin{choices}
\choice Heap sort e selection sort
\choice Quicksort e merge sort
\choice Merge sort e bubble sort
\choice Bubble sort e Quick sort
\choice Merge sort e heap sort
\end{choices}

\question
Considere as seguintes tabelas em uma base de dados relacional:\\
\textbf{Departamento (CodDepto, NomeDepto)}\\
\textbf{Empregado (CodEmp, NomeEmp, CodDepto)}\\
Deseja-se obter uma tabela na qual cada linha é a concatenação de uma linha da tabela Departamento com uma linha da tabela de Empregado. Caso um departamento não possua empregados, sua linha no resultado deve conter vazio (NULL) nos campos referentes ao empregado. A operação de álgebra relacional que deve ser aplicada para combinar estas duas tabelas é:
\begin{choices}
\choice Projeção
\choice Junção externa
\choice União
\choice Divisão
\choice Junção interna
\end{choices}

\question
Considere um sistema distribuído onde cada nó precisa obter um bloqueio (\textit{lock}) antes de acessar qualquer serviço no sistema. Qual das estratégias a seguir \textbf{não} seria eficaz para evitar impasses (\textit{deadlocks})?
\begin{choices}
\choice Instalar um serviço de detecção de impasses no sistema distribuído e reiniciar os nós que atinjam um impasse
\choice Fazer com que cada nó reinicie sua execução se um pedido de bloqueio não é concedido após um longo tempo de espera. O pedido de bloqueio é re-enviado após um tempo aleatório
\choice Numerar os serviços e exigir que cada nó solicite os bloqueios dos serviços em ordem crescente
\choice Associar prioridades aos nós e criar filas de prioridades para cada serviço
\choice Forçar cada nó a obter todos os bloqueios de que necessita no início de sua execução e reiniciar a execução se algum bloqueio não é concedido
\end{choices}

\question
Considere as afirmações abaixo sobre endereços IP:
\begin{itemize}
    \item[I.] Uma rede IP classe C fornece até 256 endereços válidos para serem atribuídos a equipe.
    \item[II.] A quantidade máxima de bits que pode ser utilizada para se definir sub-redes em uma rede IP classe C é seis (6).
    \item[III.] A máscara padrão para uma rede classe B é 255.255.255.0.
\end{itemize}
Qual das seguintes alternativas representa as afirmações corretas?
\begin{choices}
\choice Somente II.
\choice Somente II e III.
\choice Somente III.
\choice Somente I.
\choice Somente I e II.
\end{choices}

\question
A interposição de um circuito de memória cache entre o processador e a memória principal (RAM)
\begin{choices}
\choice aumenta o tráfego de instruções e/ou dados no barramento de memória
\choice permite acessos concorrentes à memória RAM
\choice diminui o tráfego de instruções e/ou dados entre memória e disco
\choice diminui o tráfego de instruções e/ou dados no barramento de memória
\choice aumenta o tráfego de instruções e/ou dados entre memória e disco
\end{choices}



    %%%%%%%%%%%%%%%%%%%%%%%%%%%%%%%%%%%%%%%%%%%%%%%%%%%%%%%%%%%%%%%%%%%%%%%%%%%%%%%%
    %%% Finaliza o bloco de questões:
    \end{questions}
    \end{document}
    